\section{Análise}

Os métodos como pode ser visto pelos resultados na tabela \ref{2_comp},
raramente detectaram os mesmos padrões. Tal é de certa forma esperado, como
apontado por pelo autor do algorítimo
probabilístico\cite{probabilistic_api_miner}, o método é feito para
especificamente detectar também padrões menos representados no dataset completo.
O método de mineração de regras de associação, por outro lado, por definição
\cite{minig_association_rules}, é criado para detectar padrões que estão
amplamente representados no dataset.

Como pode-se observar pelas tabelas \ref{desc_item}, \ref{desc_clustering} e
\ref{desc_pam}, todos os métodos detectam padrões de, em média, 2 funções. Tal é
esperado de certa forma pois, quanto mais curto, mais fácil é de achar
instâncias no dataset que suportem o padrão em questão.

Outro ponto que se aplica a todos os métodos é a relativa baixa
representatividade dos padrões detectados no dataset original. Uma possível
explicação para tal efeito é o fato que padrões diferentes se aplicam em nichos
diferentes. O padrão \textit{(pthread\_mutex\_lock,pthread\_mutex\_unlock)}
exemplifica tal fenômeno pois se aplica, quase exclusivamente, em contextos onde
threads estão em uso.

Suportando tal tese também é o fato de funções ligadas a alocação de memória e a
interação com arquivos estavam frequentemente em padrões detectados. Sendo a
gestão de memoria e o uso de arquivos algo mandatório em C, não é surpreendente
que estes padrões sejam tão sobre representados nos resultados.

Com tal visão, não é surpreendente que as \% de utilização
dos padrões detectados sejam relativamente baixas no dataset.
